\section{基于基函数的仿真求解}

\subsection{广义坐标与基函数变换}
对于不可伸长细杆而言，主导其宏观动力学行为的不是沿轴向的伸缩，而是相邻线段之间旋转状态的变化。换言之，细杆的大尺度形变主要体现为弯曲与扭转在链式结构中的传播，而不是各离散顶点在三维空间中的独立平移。因此，如果仍然把每个顶点位置与每个线段姿态都作为相互独立的原始变量，不可伸长条件所诱发的高刚度就会反复出现在底层系统中。与之相对，更自然的思路是直接围绕杆的局部旋转自由度构造广义坐标，并通过链式递推关系恢复中心线形状，从而使轴向硬约束不再作为显式变量持续出现。
材料坐标系对应一个旋转 $R\in SO(3)$，可采用多种方式表示。其李代数是一个 $3 \times 3$ 的斜对称矩阵，也可紧凑写成一个三维向量，即轴角表示。本文用 $\omega_i\in\mathcal{R}^3$ 替换第 $i$ 个线段上的四元数 $\q_i$。相应的单位四元数可由 $\omega_i$ 直接计算得到：
\begin{equation}
  \q_i = \mathcal{B}\left(\omega_i\right) = \left[\cos \frac{\left\|
      \omega_i \right\|}{2}, \frac{\omega_i}{\left\| \omega_i \right\| }
    \sin \frac{\left\| \omega_i \right\|}{2}\right]^\top.
    \label{eq: omega2quaternion}
\end{equation}
当 $\left\| \omega_i \right\|$ 趋近于零时，应用洛必达法则来避免除零问题：
\begin{equation}
  \q_i = \Lim{\left\| \omega_i \right\| \rightarrow
    0} \mathcal{B}\left(\omega_i\right)
    = \left[1, \frac{\omega_i}{2} \right]^\top.
    \label{eq: omega2quaternion2}
\end{equation}
这一选择消除了 \Cref{eq: unitq_cons} 中的单位四元数约束。

{\color{black}相应的 $3\times 3$ 旋转矩阵也可依据罗德里格斯公式，通过指数映射写为：}
\begin{equation}
\mathbf{m}_i
= \operatorname{exp}\left( \mathbf{\omega}_i\right).
\end{equation}
采用图~\ref{fig:single-chain} 所示的链式表示后，不可伸长杆上非根顶点的位置可递归写为：
\begin{equation}
\p_{i+1}=\p_i+l_id(\mathbf{m}_i)=\p_i+l_id(\operatorname{exp}(\omega_i)), i=1,2,\cdots,n-1.
\label{eq:w2p}
\end{equation}
这同时消除了 \Cref{eq: alignment_cons} 中的对齐约束与 \Cref{eq: inext_cons_d} 中的不可伸长约束。线段对齐轴也可由四元数形式 $d(\mathcal{B}(\omega_i))$ 得到，但计算代价更高。

从运动学角度看，这种链式递推关系的意义在于：每个线段的位置不再是需要独立求解的自由变量，而是由根节点位置、线段静止长度以及沿链条传播的局部旋转共同决定。因此，不可伸长条件不再需要通过额外的距离方程反复维持，而是由位置更新规则本身保证。换言之，链式表示并不是简单地重新标记原有变量，而是通过改变自由度的组织方式，从变量层面移除最显式的拉伸硬约束。

最后，本文用笛卡尔坐标 $\p_c=\p_0$ 表示根顶点位置。

\begin{figure}[H]
  \centering
  \centering
  \includegraphics[width=0.75\textwidth]{single_chain_new.pdf}
  \caption[链式结构]{本文采用一种不含冗余自由度的紧凑链式结构来表示不可伸长柯塞拉特杆，即基尔霍夫杆。}
  \label{fig:single-chain}
\end{figure}

综合以上内容，本文使用根顶点处的 $\p_c \in \mathcal{R}^3$ 和第 $i$ 个线段 $\mathbf{e}_i$ 上的 $\omega_i\in\mathcal{R}^3$ 来表示整根杆。传统表示使用 $\p$ 和 $\q$ 中的 $7n-4$ 个变量，并受到~\Cref{eq: inext_cons_d,eq: unitq_cons,eq: alignment_cons} 中 $4(n-1)$ 个线性独立约束限制，因此实际自由度为 $(7n-4)-4(n-1)=3n$。本文的表示方法 $\mathbf{s}=[\p_{\text{c}}^\top, \mathbf{\omega}^\top]^\top$ 恰好也有 $3n$ 个自由度，与非闭环不可伸长分段线性柯塞拉特杆的自由度一致。采用该表示后，\Cref{eq: inext_cons_d,eq: unitq_cons,eq: alignment_cons} 中的等式约束均由参数化自动满足。\textcolor{black}{实际上，该方法有效地建模了一类特殊的不可剪切且不可伸长的柯塞拉特杆，即基尔霍夫杆~\cite{dill1992kirchhoff}。}

这一点也说明了本文方法的物理适用范围。通过将位置更新完全建立在线段静止长度与局部旋转之上，本文事实上默认了轴向拉伸与剪切形变可以忽略，而主要保留弯曲与扭转主导的动力学响应。对于细长比很大、且几何变化主要体现为弯扭传播的线缆、细丝与线状结构，这一假设是合理的；也正是在这种场景下，采用更贴近杆内在运动学的广义坐标，能够显著减少传统表示中的冗余约束负担。

在传统表示方法下，隐式欧拉时间步进会导致如下非线性约束优化问题~\cite{martin2011example}：
\begin{subequations}
  \begin{align}
    &\mathbf{x}_{t+\Delta t} = \mathop{\arg\min}\limits_{\mathbf x} E(\mathbf{x}) =
  V\left(\mathbf{q}\right)+T_p\left(\mathbf{p}\right)+T_q\left(\mathbf q \right)+G\left(\mathbf{p}\right),\\
    &\text{s.t.}  \left\{
  \begin{array}{l}
  \mathbf{C_E}(\mathbf{x}) = \mathbf{0},\\
    \mathbf{C_I} (\mathbf{x})\geq \mathbf{0}.
    \label{eq:cons_origin_opt}
  \end{array}
      \right.
     % \quad \mathbf{C_E}(\mathbf{x}) = \mathbf{0}\label{eq: eqcons}\\
    %&\quad \quad \mathbf{C_I} (\mathbf{x})\geq \mathbf{0} \label{eq: ineqcons}
  \end{align}
  \label{eq: origin_opt}
\end{subequations}

\vspace{-0.5cm}
上述方程中，$V(\mathbf q)$ 为离散弯曲和扭转势能，$T_p(\mathbf p)$ 和 $T_q(\mathbf q)$ 分别为平动和转动动能，$G\left(\mathbf{p}\right)$ 为外力场（如重力）的势能。这些项的具体公式参见\ref{app:energies}节，更多细节可参阅文献~\cite{spillmann_corde_2007,wen_cosserat_2020}。

\inputbody{paper1/app_energy_terms}
\Cref{eq:cons_origin_opt} 中的大量等式约束 $\mathbf{C_E}$ 来自~\cref{eq: inext_cons_d,eq: unitq_cons,eq: alignment_cons}。由~\cref{eq: omega2quaternion,eq: omega2quaternion2,eq:w2p} 和 $\p_c=\p_0$ 可导出变换 $\mathcal{T}$，从而直接用紧凑变量 $\mathbf{s}$ 参数化满足几何约束的杆构型：
\begin{equation}
  \mathbf{x} = \mathcal{T}\left( \mathbf{s} \right).
\end{equation}

换言之，优化变量已经从带有显式距离、单位长度和对齐约束的传统空间，转移到天然满足这些约束的变换后坐标空间 $\mathbf{s}$ 中。因此，新表示下的隐式时间步进可重写为一个仅需处理剩余边界条件与接触约束的优化问题。

需要指出的是，约束的减少并不意味着问题完全失去了非线性。经过表示重构后，细杆的内在几何约束已经不再显式出现，但外部边界条件、狄利克雷约束以及接触所诱发的不等式条件仍然需要在每个时间步中得到一致处理。因此，表示重构的作用并不是把原问题变成一个无约束的简单优化，而是将求解器所面对的约束集合压缩到真正与外部环境和边界交互相关的部分。

本文采用经典的主动集框架~\cite{nocedal1999numerical} 来处理剩余的等式约束（如狄利克雷边界条件约束）和通常来自碰撞的不等式约束 $\mathbf{C_I}$。具体而言，在主动集法的每次迭代中，求解如下问题：
\begin{equation}
\begin{gathered}
    {\mathbf{s}}_{t+\Delta t} = \mathop{\arg \min}\limits_{{\mathbf s}} E(\mathcal{T}\left({\mathbf{s}}\right)), \\
    \text{s.t.}\quad
    \mathbf{C_{\mathcal A}} (\mathcal{T}\left({\mathbf{s}} \right))= \mathbf{0}.
    \end{gathered}
    \label{eq: trans_opt}
\end{equation}
约束 $\mathbf{C_{\mathcal A}}$ 会根据当前碰撞状态和剩余等式约束动态更新。具体而言，剩余等式约束始终保留在集合 $\mathbf{C_{\mathcal A}}$ 中；$\mathbf{C_I}$ 中处于活动状态的不等式约束（如活动碰撞）会被加入该集合，并按等式约束处理；非活动碰撞则在每次迭代中移除。每一步迭代后，系统状态都可回溯到非穿透构型。伪代码见附录\ref{app:active_set}。

通过上述轴角链式广义坐标变换，细杆的内在极硬约束已在表示层面消除。此时，系统的隐式时间积分转化为一个维数更低、且仅受碰撞等外部边界条件限制的优化问题。为了高效处理这些剩余的不等式约束并保持大时间步长下的稳定性，本文在新表示下引入基于主动集的序列二次规划（SQP）求解框架。
\subsection{新表示下的 SQP 求解框架}
总体上，本文遵循 Pan 等人~\cite{pan_subspace_2015} 的思路，在序列二次规划（Sequential Quadratic Programming, SQP）框架下（算法~\ref{alg: nonlinear_solver}）求解~\cref{eq: trans_opt} 中的非线性优化问题。为提高 SQP 的收敛性，本文采用线搜索策略~\cite{nocedal1999numerical}，并用下式价值函数评估步长：
\begin{equation}
   \phi(\mathbf{s}, \mu) = E(\mathcal{T}(\mathbf{s})) + \frac{1}{\mu}\left\|\mathbf{C}_{\mathcal{A}} \right\|_{1}.
   \label{eq:merit_function}
\end{equation}
\textcolor{black}{
为防止步长过度违反约束，$\mu$ 被更新为当前活动约束对应拉格朗日乘子最大值的倒数~\cite{nocedal1999numerical}}。

具体地，每一步在当前 $\mathbf s^k$ 处构造~\cref{eq: trans_opt} 的二次模型如下：
\begin{equation}
\begin{gathered}
    \Delta {\mathbf{s}}^k =\mathop{\arg \min}\limits_{\Delta \mathbf s}
    \frac{1}{2}\Delta{\mathbf{s}}^{\top} \mathbf{A}^k \Delta {\mathbf{s}} +
    \mathbf{b}^{k^\top} \Delta {\mathbf{s}},\\
    \text{s.t.} \quad
    \nabla_{{\mathbf{s}}}\mathbf{C_{\mathcal A}}  \Delta{\mathbf{s}} + \mathbf{C_{\mathcal A}}(\mathcal{T}({\mathbf{s}}^k )) = \mathbf{0},
  \label{eq: quad_opt}
  \end{gathered}
\end{equation}
其中，设 $\J(\mathbf{s}) = D \mathcal{T} = \partial \mathbf{x}/\partial \mathbf{s}$，有
\begin{equation*}
\begin{split}
\mathbf{b}^k &= \J(\mathbf{s}^k)^\top \left
  .\FPP{E(\mathbf{x})}{\mathbf{x}}\right|_{\mathbf{x} =
  \mathcal{T}(\mathbf{s}^k)},\\
\mathbf{A}^k &= \J(\mathbf{s}^k)^\top \left.
  \FPP{^2E(\mathbf{x})}{\mathbf{x}^2}\right|_{\mathbf{x} = \mathcal{T}(\mathbf{s}^k)}
\J(\mathbf{s}^k) = \J^{k\top} \mathbf{H}^k \J^k.
\end{split}
\end{equation*}

这里，本文与 Pan 等人~\cite{pan_subspace_2015}一样，省略 $\left.\FPP{\J\left(\mathbf{s}\right)}{\mathbf{s}}\right|_{\mathbf s=\mathbf{s}^k}:\left.\FPP{E(\mathbf{x})} {\mathbf{x}}\right|_{\mathbf{x} = \mathcal{T}\left(\mathbf{s}^k\right)}$。因此，每个二次子问题实际上采用不精确牛顿法求解。

这样处理主要是为了平衡计算代价与收敛质量。对于链式表示而言，若继续保留由 $\J(\mathbf{s})$ 的导数产生的完整二阶项，则需要显式处理更复杂的高阶张量收缩，其代价会明显削弱紧凑表示在算子结构上的优势。相较之下，不精确牛顿法虽然牺牲了部分局部二阶精度，但仍保留了主要曲率信息，并能在每次迭代中以更低成本构造可解的二次子问题。对本文关注的大规模细杆隐式仿真而言，这种取舍更符合整体效率目标。
\begin{algorithm}
\caption{SQP 框架}\label{alg: nonlinear_solver}
\begin{algorithmic}
\REQUIRE $\mathbf{x}^k$
\STATE $\mathbf{s}^k = \mathcal{T}^{-1}(\mathbf{x}^k)$
\WHILE{$\big\| \mathbf b + \nabla_{\mathbf{s}} \mathbf{C}_{\mathcal A}^\top \mathbf\lambda \big\|_{2} \neq 0$}
   \STATE 评估 $\mathbf{A}^k$，$\mathbf{b}^k$，$\nabla_{\mathbf{s}}\mathbf{C}_{\mathcal A}^k$
   \STATE $\Delta {\mathbf{s}}$, $\lambda$ = active set( $\mathbf{A}^k$, $\mathbf{b}^k$, $\nabla_{\mathbf{s}}\mathbf{C}_{\mathcal A}^k$)
   \STATE $\alpha$ = line search ($\mathbf{s}^k$, $\Delta \mathbf{s}$, $\lambda$)
   \STATE ${\mathbf{s}}^{k+1} = {\mathbf{s}}^k + \alpha \Delta {\mathbf{s}}$
\ENDWHILE
\end{algorithmic}
\end{algorithm}

%We should also notice that the $\nabla_{\mathbf{s}}\mathbf{C_E}$ and
%$\nabla_{\mathbf{s}}\mathbf{C_E}$ are both stay unchanged. {\color{black}Also we
%  should notice that our hessian $\mathbf{A}$ may be very dense but easy to
%  employ matrix-vector product and inverse-vector operation. which asking the
%  solver to be matrix-free.}
\Cref{eq: quad_opt} 中的 QP 问题导致一个 KKT 系统：
\begin{equation}
  \begin{bmatrix}
     \mathbf{A}  &  \nabla_{{\mathbf{s}}} \mathbf{C}_{\mathcal{A}}^\top
    \\\nabla_{{\mathbf{s}}} \mathbf{C}_{\mathcal{A}}   & \bold{0}
  \end{bmatrix}
  \begin{bmatrix} \Delta {\mathbf{s}} \\ \lambda \end{bmatrix}
  =\begin{bmatrix}
  \mathbf{b}\\
    - \mathbf{C}_{\mathcal{A}}(\mathcal{T}\left({\mathbf{s}}_k)\right)
  \end{bmatrix}.
  \label{eq: final_kkt}
\end{equation}
%where $\mathbf{C} = \left[\mathbf{C}_{\text{active}}^\top \right]^\top$.
通过\textit{舒尔补}方法，先求解 $\lambda$ 然后求解 $\Delta {\mathbf{s}}$：
% where in this KKT the constraints $C$ include the active set of $\mathbf{C_I}$.
% This KKT is solved by ranged-based method:
\begin{equation}
\left\{
\begin{array}{l}
%\begin{split}
\left( \nabla_{\mathbf{s}}\mathbf{C}_{\mathcal{A}}\mathbf{A}^{-1} \nabla_{\mathbf{s}} \mathbf{C}_{\mathcal{A}}^{\top}\right)  \lambda =\mathbf{C}_{\mathcal{A}} -  \nabla_{\mathbf{s}}\mathbf{C}_{\mathcal{A}}\mathbf{A}^{-1} \mathbf{b},\\
\mathbf{A}\Delta {\mathbf s} =  \mathbf{b} - \nabla_{\mathbf{s}} \mathbf{C}_{\mathcal{A}}^\top \lambda.
% \end{split}
\end{array}
\right.
\label{eq: solve_kkt}
\end{equation}

与 RedMax~\cite{wang_redmax_2019} 类似，本文使用预条件共轭梯度（Preconditioned Conjugate Gradient, PCG）方法求解 $\mathbf{A}$ 对应的方程。这类方程的数量等于 $\nabla_s\mathbf{C_{\mathcal A}}$ 的行数加上 $2$，会带来可观的计算成本。PCG 求解线性系统的成本与矩阵--向量乘法及预条件子密切相关；若充分利用 $\mathbf{A}$ 的结构，两类操作的时间复杂度均可达到 $\mathcal{O}(n)$，同时还能加速收敛。详细分析见下一节。

\inputbody{paper1/app_active_set}

至此，表示重构已经将最显式的几何硬约束从变量层面移除，但数值求解的主要瓶颈并未完全消失。每一次 SQP 迭代仍然需要反复评估雅可比、黑塞矩阵以及由此产生的线性系统；若这些关键算子的计算与求解仍按稠密矩阵处理，前面通过坐标重构获得的结构优势就难以真正转化为计算效率。因此，下一节的重点不再是重新讨论约束如何消除，而是说明重构后的系统矩阵为何具有可利用的结构，以及这种结构如何进一步支持线性复杂度的算子评估与预条件子设计。

为了在迭代求解中维持较高效率，需要降低雅可比和黑塞矩阵相关关键算子的计算复杂度。以下将详细阐述其线性时间复杂度的构造过程。
\subsection{\texorpdfstring{$\mathcal{O}(n)$}{O(n)} 时间复杂度的算子评估}
\label{ssec:jacobian-patt}

PCG 算法中需要频繁对矩阵 $\mathbf{A} = \J^\top\mathbf{H}\J$ 执行矩阵--向量乘法。下文将证明，在紧凑表示下，将 $\mathbf{H}, \J, \J^\top, \mathbf{A}$ 与向量相乘都可以在 $\mathcal{O}(n)$ 时间内实现。

为简化表述，采用如下记号：
\begin{itemize}
  \item $\LowerTri$：块元素均为 $\I$ 的块下三角矩阵，
  \item $\UpperTri$：块元素均为 $\I$ 的块上三角矩阵，
  \item $\Diag$：块对角矩阵（忽略其具体数值），
  \item $\TriDiag$：块三对角矩阵（忽略其具体数值）。
\end{itemize}

\textbf{$\mathbf H$ 的模式：}对于杆仿真，在应用变换 $\mathcal T$ 之前，系统矩阵具有简单的带状结构，采用 $(\mathbf{p}, \mathbf{q})$ 排序时如下：
\begin{equation}
  \label{eq:H-patt}
  \mathbf{H} =
  \begin{bmatrix}
    \mathbf{H}_{\mathbf{p}} & \mathbf{0}              \\
    \mathbf{0}              & \mathbf{H}_{\mathbf{q}}
  \end{bmatrix} =
  \begin{bmatrix}
    \frac{1}{\Delta t^{2}}\mathbf{M}_{\mathbf{p}} & \mathbf{0}                                                               \\
    \mathbf{0}                                    & \frac{1}{\Delta t^{2}} \mathbf{M}_{\mathbf{q}} + \mathbf{K}_{\mathbf{q}}
  \end{bmatrix} =
  \left[
    \begin{array}{@{}c|c@{}}
      \Diag      & \mathbf{0} \\
      \hline
      \mathbf{0} & \TriDiag
    \end{array}
    \right],
\end{equation}
其中 $ \mathbf{H}_{\p} = \FPP{^2E}{\mathbf{p}^2} = \FPP{^2T_{p}}{\mathbf{p}^2} =\mathbf{M}_\p/\Delta t^2$，且 $\mathbf{H}_{\q}= \FPP{^2E}{\mathbf{q}^2} = \FPP{^2T_{q}(\mathbf{q})}{\mathbf{q}^2} + \FPP{^2V(\mathbf{q})}{\mathbf{q}^2} = \mathbf{M}_q /\Delta t^2 + \mathbf{K}_q$。在这种简单的带状结构下，$\mathbf H$ 的矩阵 - 向量乘法的 $\mathcal{O}(n)$ 时间复杂度自然得到保证。

\textbf{$\mathbf{J}$ 的模式：}将雅可比矩阵 $\J$ 分为两块：
\begin{equation}
  \J =
  \begin{bmatrix}
    \J_{\mathbf{p, s}} \\
    \J_{\mathbf{q, s}}
  \end{bmatrix}.
\end{equation}
由于 $\mathbf s=[\p_{\text{c}}^\top, \mathbf{\omega}^\top]^\top$，$\J_{\mathbf{q, s}}$ 具有如下模式：
\begin{equation}
  \J_{\mathbf{q, s}} =
  \begin{bmatrix}
    \frac{\partial \mathbf{q}}{\partial \mathbf{p}_{c}} & \frac{\partial \mathbf{q}}{\partial \mathbf{\omega}}
  \end{bmatrix} =
  \begin{bmatrix}
    \mathbf{0} & \Diag
  \end{bmatrix}.
  \label{eq:Jqs_pattern}
\end{equation}
$\J_{\mathbf{p, s}}$ 的模式相对复杂，如下所示：
\begin{equation}
  \begin{aligned}
    \J_{\mathbf{p, s}} & =
    \begin{bmatrix}
      \I     &                        &                                                 &                        \\
      \I     & \J^{1}_{\mathbf{p, s}} & \multicolumn{2}{c}{\raisebox{1ex}[0pt]{\Huge0}}                          \\
      \vdots & \vdots                 & \ddots                                          &                        \\
      \I     & \J^{1}_{\mathbf{p, s}} & \cdots                                          & \J^{n}_{\mathbf{p, s}} \\
    \end{bmatrix} \\
                       & =
    \begin{bmatrix}
      \I     &        &                                                      \\
      \I     & \I     & \multicolumn{2}{c}{\raisebox{1ex}[0pt]{\Huge0}}      \\
      \vdots & \vdots & \ddots                                          &    \\
      \I     & \I     & \cdots                                          & \I \\
    \end{bmatrix}
    \underbrace{
    \begin{bmatrix}
        \I &                        &        &                        \\
           & \J^{1}_{\mathbf{p, s}} &        &                        \\
           &                        & \ddots &                        \\
           &                        &        & \J^{n}_{\mathbf{p, s}}
      \end{bmatrix}}_{\mathrm{Diag}\left(\{\J^{i}_{\mathbf{p, s}}\}\right)}                                      \\&=
    \begin{bmatrix}
      \LowerTri \cdot \Diag
    \end{bmatrix}.
  \end{aligned}
\end{equation}
由上述结构可知，$\J_{\mathbf{q, s}}$ 与向量的乘法可在 $\mathcal{O}(n)$ 时间内完成。$\J_{\mathbf{p, s}}$ 的矩阵--向量乘法复杂度与 $\LowerTri$ 相同，相当于在 $\mathcal{O}(n)$ 时间内计算向量的前缀和~\cite{harris2007parallel}。因此，$\mathbf J$ 和 $\J^{\top}$ 的矩阵--向量乘法均可在线性时间内实现。

综上，$\mathbf{A}$ 与向量的乘法也可在 $\mathcal{O}(n)$ 时间内完成：
\begin{equation}
  \mathbf{A} = \J^{\top}\mathbf{H}\J =
  \underbrace{
    \left[
      \begin{array}{@{}c|c@{}}
        \Diag \cdot \UpperTri &
        \begin{bmatrix}
          \mathbf{0} \\
          \Diag
        \end{bmatrix}
      \end{array}
    \right]}_{\J^{\top}}
  \underbrace{
    \left[
      \begin{array}{@{}c|c@{}}
        \Diag      & \mathbf{0} \\
        \hline
        \mathbf{0} & \TriDiag
      \end{array}
    \right]}_{\mathbf{H}}
  \underbrace{
    \left[
      \begin{array}{c}
        \LowerTri \cdot \Diag \\
        \hline
        \begin{bmatrix}
          \mathbf{0} & \Diag \\
        \end{bmatrix}
      \end{array}
    \right]}_{\J}.
\end{equation}
\inputbody{paper1/preconditioner_content}

%\textcolor{black}{As a result, the performance can be improved $10$ to $100$ times.}
